% !TEX root = ../../Main.tex
\section{Limitations}
\label{Section:Limitations}

The most important limitation comes first, because it puts a limit on everything else. The performance figures in \Cref{Chapter:ExperimentalResults} are measured over a window that includes the years the agents were trained on. The procedure does keep one year out of training and out of the choice among an agent's own episodes. But the choice of which trained agent to publish was made on the full window, so the headline returns are an in-sample result. On the one year that no part of the procedure saw, the models lost money. The portfolio returned $-9.8\%$ in 2025, and five of the seven pairs were negative. That year had a broad dollar reversal, and the training decade had none. The one model that held up is the least directional model in the set. However, a single year is a small sample, and no stronger conclusion should be drawn from it in either direction.

There is a second way in which the selection is not corrected. Each published model was chosen from between $128$ and $256$ candidates. With that many draws, a permutation $p$-value near $0.05$ is close to what chance alone produces. For this reason, the screening is reported as a filter that was shown to reject bad models. It is not reported as a claim that the models that passed are statistically significant. Under a strict correction, only USD/CAD and GBP/USD would clearly pass.

Three smaller limitations should also be recorded. First, the position size was chosen for each pair after the results were seen. It scales every trade and does not change the policy that makes the trades. Still, it is a parameter fitted on the same data as everything else. Second, training cannot be reproduced across environments. It is bit-exact within one installation on a single thread, but a different version of a numerical library changes the outcome. So the delivered models are archived artifacts whose replay has been verified. They are not a recipe that can produce the same models again. Third, two weeks of January 2016 are missing from all seven pairs, for reasons on the broker's side. The gap is the same for all pairs and lies outside the evaluation window, so it does not bias any comparison in this work.

Finally, the tests cover a narrower scope than the method itself. They use seven instruments from one asset class, one eleven-year period, one decision frequency and one cost model. This is only one case out of many possible ones. Nothing in this work shows that the design carries over to other markets, other periods or other holding horizons. The USD/JPY result shows that it does not carry over to every instrument, even among the seven.